\section{为什么特别强调效率}

课程把效率定义成一个总括性的概念：数据、算力、内存、带宽和算法共同决定模型能做到什么。
这也解释了为什么课程安排不是按“模型论文”来排，而是按“效率瓶颈”来排。

\begin{importantbox}{效率驱动设计}
\begin{itemize}
    \item 数据处理要避免把计算浪费在低质量数据上。
    \item Tokenization 要在表达能力和计算效率之间折中。
    \item 模型结构要考虑内存、FLOPs 和硬件友好性。
    \item 训练要和预算匹配，不是越久越好。
    \item 对齐和数据选择也会回到资源效率问题。
\end{itemize}
\end{importantbox}

